\chapter{Objekte, Spieler und Entities}
\label{ch:map/objects}

In diesem Kapitel lernst du, wie du Objekte veränderst und eigene
definierst.

\section{Was sind Objekte, Spieler und
Entities?}\label{was-sind-objekte-spieler-und-entities}

Spieler und Entities\index{entity@Entity} sind beides Objekte. Ein
Objekt ist etwas, das sich unabhängig vom Node-Gitter bewegen kann und
Eigenschaften wie Geschwindigkeit und Größe hat. Objekte sind keine
Items und haben ein eigenes Registrierungssystem.

Zwischen Spielern und Entities gibt es einige Unterschiede. Der
wichtigste: Spieler werden von Menschen gesteuert, Entities von Mods.
Deshalb können Mods die Geschwindigkeit eines Spielers nicht setzen --
Spieler werden auf dem Client berechnet, Entities auf dem Server. Ein
weiterer Unterschied: Spieler sorgen dafür, dass Mapblöcke geladen
werden, Entities werden dagegen einfach gespeichert und inaktiv.

Verwischt wird diese Unterscheidung dadurch, dass Entities über eine
Tabelle gesteuert werden, die \emph{Lua-Entity} heißt; dazu später mehr.

\section{Position und
Geschwindigkeit}\label{position-und-geschwindigkeit}

Mit \texttt{get\_\allowbreak{}pos} und \texttt{set\_\allowbreak{}pos}
liest und setzt du die Position eines Objekts.

\begin{Shaded}
\begin{Highlighting}[]
\KeywordTok{local} \VariableTok{object} \OperatorTok{=} \VariableTok{core}\OperatorTok{.}\NormalTok{get\_player\_by\_name}\OperatorTok{(}\StringTok{"bob"}\OperatorTok{)}
\KeywordTok{local} \VariableTok{pos} \OperatorTok{=} \VariableTok{object}\OperatorTok{:}\NormalTok{get\_pos}\OperatorTok{()}
\VariableTok{object}\OperatorTok{:}\NormalTok{set\_pos}\OperatorTok{(\{} \VariableTok{x} \OperatorTok{=} \VariableTok{pos}\OperatorTok{.}\VariableTok{x}\OperatorTok{,} \VariableTok{y} \OperatorTok{=} \VariableTok{pos}\OperatorTok{.}\VariableTok{y} \OperatorTok{+} \DecValTok{1}\OperatorTok{,} \VariableTok{z} \OperatorTok{=} \VariableTok{pos}\OperatorTok{.}\VariableTok{z} \OperatorTok{\})}
\end{Highlighting}
\end{Shaded}

\texttt{set\_\allowbreak{}pos} setzt die Position sofort, ohne
Animation. Soll das Objekt sich fließend an die neue Position bewegen,
verwendest du \texttt{move\_\allowbreak{}to}. Das funktioniert leider
nur bei Entities.

\begin{Shaded}
\begin{Highlighting}[]
\VariableTok{object}\OperatorTok{:}\NormalTok{move\_to}\OperatorTok{(\{} \VariableTok{x} \OperatorTok{=} \VariableTok{pos}\OperatorTok{.}\VariableTok{x}\OperatorTok{,} \VariableTok{y} \OperatorTok{=} \VariableTok{pos}\OperatorTok{.}\VariableTok{y} \OperatorTok{+} \DecValTok{1}\OperatorTok{,} \VariableTok{z} \OperatorTok{=} \VariableTok{pos}\OperatorTok{.}\VariableTok{z} \OperatorTok{\})}
\end{Highlighting}
\end{Shaded}

Bei Entities musst du immer an die Netzwerklatenz denken. Im Idealfall
kämen Nachrichten über Bewegungen sofort an, in der richtigen
Reihenfolge und in ähnlichen Abständen, wie du sie gesendet hast.
Außerhalb des Einzelspielermodus ist die Welt aber nicht ideal.
Nachrichten brauchen eine Weile. Positionsnachrichten können in falscher
Reihenfolge ankommen, sodass manche
\texttt{set\_\allowbreak{}pos}-Aufrufe übersprungen werden -- es hat
keinen Sinn, an eine Position zu springen, die älter ist als die aktuell
bekannte. Bewegungen kommen womöglich nicht in gleichmäßigen Abständen
an, was sie für Animationen schwer nutzbar macht. All das führt dazu,
dass der Client etwas anderes sieht als der Server. Dessen solltest du
dir bewusst sein.

\section{Objekteigenschaften}\label{objekteigenschaften}

Objekteigenschaften sagen dem Client, wie er ein Objekt darstellen und
behandeln soll. Eigene Eigenschaften kannst du nicht festlegen, denn die
Eigenschaften sind per Definition für die Engine bestimmt.

Anders als Nodes haben Objekte kein festes, sondern ein veränderliches
Aussehen. Wie ein Objekt aussieht -- und einiges mehr --, kannst du
jederzeit ändern, indem du seine Eigenschaften aktualisierst.

\begin{Shaded}
\begin{Highlighting}[]
\VariableTok{object}\OperatorTok{:}\NormalTok{set\_properties}\OperatorTok{(\{}
    \VariableTok{visual} \OperatorTok{=} \StringTok{"mesh"}\OperatorTok{,}
    \VariableTok{mesh} \OperatorTok{=} \StringTok{"character.b3d"}\OperatorTok{,}
    \VariableTok{textures} \OperatorTok{=} \OperatorTok{\{} \StringTok{"character\_texture.png"} \OperatorTok{\},}
    \VariableTok{visual\_size} \OperatorTok{=} \OperatorTok{\{} \VariableTok{x} \OperatorTok{=} \DecValTok{1}\OperatorTok{,} \VariableTok{y} \OperatorTok{=} \DecValTok{1} \OperatorTok{\},}
\OperatorTok{\})}
\end{Highlighting}
\end{Shaded}

Die geänderten Eigenschaften werden an alle Spieler in Reichweite
geschickt. So erreichst du sehr günstig viel Abwechslung, etwa
unterschiedliche Skins für jeden Spieler.

Wie der nächste Abschnitt zeigt, können Entities in ihrer Definition
Anfangseigenschaften haben. Die Standardeigenschaften von Spielern sind
dagegen in der Engine festgelegt; für neu beigetretene Spieler setzt du
die Eigenschaften deshalb mit \texttt{set\_\allowbreak{}properties()} in
\texttt{on\_\allowbreak{}joinplayer}\index{callback@Callback!on_joinplayer@\texttt{on\_\allowbreak{}joinplayer}}.

\section{Entities}\label{entities}

Eine Entity hat eine Definitionstabelle, die einer Item-Definition
ähnelt. Sie kann Callback-Methoden, Anfangseigenschaften des Objekts und
eigene Felder enthalten.

\begin{Shaded}
\begin{Highlighting}[]
\KeywordTok{local} \VariableTok{MyEntity} \OperatorTok{=} \OperatorTok{\{}
    \VariableTok{initial\_properties} \OperatorTok{=} \OperatorTok{\{}
        \VariableTok{hp\_max} \OperatorTok{=} \DecValTok{1}\OperatorTok{,}
        \VariableTok{physical} \OperatorTok{=} \KeywordTok{true}\OperatorTok{,}
        \VariableTok{collide\_with\_objects} \OperatorTok{=} \KeywordTok{false}\OperatorTok{,}
        \VariableTok{collisionbox} \OperatorTok{=} \OperatorTok{\{} \OperatorTok{{-}}\DecValTok{0.3}\OperatorTok{,} \OperatorTok{{-}}\DecValTok{0.3}\OperatorTok{,} \OperatorTok{{-}}\DecValTok{0.3}\OperatorTok{,} \DecValTok{0.3}\OperatorTok{,} \DecValTok{0.3}\OperatorTok{,} \DecValTok{0.3} \OperatorTok{\},}
        \VariableTok{visual} \OperatorTok{=} \StringTok{"wielditem"}\OperatorTok{,}
        \VariableTok{visual\_size} \OperatorTok{=} \OperatorTok{\{} \VariableTok{x} \OperatorTok{=} \DecValTok{0.4}\OperatorTok{,} \VariableTok{y} \OperatorTok{=} \DecValTok{0.4} \OperatorTok{\},}
        \VariableTok{textures} \OperatorTok{=} \OperatorTok{\{} \StringTok{""} \OperatorTok{\},}
        \VariableTok{spritediv} \OperatorTok{=} \OperatorTok{\{} \VariableTok{x} \OperatorTok{=} \DecValTok{1}\OperatorTok{,} \VariableTok{y} \OperatorTok{=} \DecValTok{1} \OperatorTok{\},}
        \VariableTok{initial\_sprite\_basepos} \OperatorTok{=} \OperatorTok{\{} \VariableTok{x} \OperatorTok{=} \DecValTok{0}\OperatorTok{,} \VariableTok{y} \OperatorTok{=} \DecValTok{0} \OperatorTok{\},}
    \OperatorTok{\},}

    \VariableTok{message} \OperatorTok{=} \StringTok{"Standardnachricht"}\OperatorTok{,}
\OperatorTok{\}}

\KeywordTok{function} \VariableTok{MyEntity}\OperatorTok{:}\NormalTok{set\_message}\OperatorTok{(}\VariableTok{msg}\OperatorTok{)}
    \VariableTok{self}\OperatorTok{.}\VariableTok{message} \OperatorTok{=} \VariableTok{msg}
\KeywordTok{end}
\end{Highlighting}
\end{Shaded}

In einem sehr wichtigen Punkt unterscheiden sich Entity-Definitionen von
Item-Definitionen: Wird eine Entity \emph{emerged} (also geladen oder
erzeugt), wird für sie eine neue Tabelle angelegt, die von der
Definitionstabelle \emph{erbt}.

Über das ObjectRef\index{objectref@\texttt{ObjectRef}} kommst du an die
Entity-Tabelle und umgekehrt:

\begin{Shaded}
\begin{Highlighting}[]
\KeywordTok{local} \VariableTok{entity} \OperatorTok{=} \VariableTok{object}\OperatorTok{:}\NormalTok{get\_luaentity}\OperatorTok{()}
\KeywordTok{local} \VariableTok{object} \OperatorTok{=} \VariableTok{entity}\OperatorTok{.}\VariableTok{object}
\FunctionTok{print}\OperatorTok{(}
    \StringTok{"Entity ist bei "}
        \OperatorTok{..} \VariableTok{core}\OperatorTok{.}\NormalTok{pos\_to\_string}\OperatorTok{(}\VariableTok{object}\OperatorTok{:}\NormalTok{get\_pos}\OperatorTok{())}
\OperatorTok{)}
\end{Highlighting}
\end{Shaded}

Für Entities gibt es eine Reihe von Callbacks.\index{callback@Callback}
Eine vollständige Liste findest du in der Lua-API-Referenz\footnote{\url{https://api.luanti.org/minetest-namespace-reference/\#registered-definition-tables}}.

\begin{Shaded}
\begin{Highlighting}[]
\KeywordTok{function} \VariableTok{MyEntity}\OperatorTok{:}\NormalTok{on\_step}\OperatorTok{(}\VariableTok{dtime}\OperatorTok{)}
    \KeywordTok{local} \VariableTok{pos} \OperatorTok{=} \VariableTok{self}\OperatorTok{.}\VariableTok{object}\OperatorTok{:}\NormalTok{get\_pos}\OperatorTok{()}
    \KeywordTok{local} \VariableTok{pos\_down} \OperatorTok{=}
        \VariableTok{vector}\OperatorTok{.}\NormalTok{subtract}\OperatorTok{(}\VariableTok{pos}\OperatorTok{,} \VariableTok{vector}\OperatorTok{.}\NormalTok{new}\OperatorTok{(}\DecValTok{0}\OperatorTok{,} \DecValTok{1}\OperatorTok{,} \DecValTok{0}\OperatorTok{))}

    \KeywordTok{local} \VariableTok{delta}
    \ControlFlowTok{if} \VariableTok{core}\OperatorTok{.}\NormalTok{get\_node}\OperatorTok{(}\VariableTok{pos\_down}\OperatorTok{).}\VariableTok{name} \OperatorTok{==} \StringTok{"air"} \ControlFlowTok{then}
        \VariableTok{delta} \OperatorTok{=} \VariableTok{vector}\OperatorTok{.}\NormalTok{new}\OperatorTok{(}\DecValTok{0}\OperatorTok{,} \OperatorTok{{-}}\DecValTok{1}\OperatorTok{,} \DecValTok{0}\OperatorTok{)}
    \ControlFlowTok{elseif} \VariableTok{core}\OperatorTok{.}\NormalTok{get\_node}\OperatorTok{(}\VariableTok{pos}\OperatorTok{).}\VariableTok{name} \OperatorTok{==} \StringTok{"air"} \ControlFlowTok{then}
        \VariableTok{delta} \OperatorTok{=} \VariableTok{vector}\OperatorTok{.}\NormalTok{new}\OperatorTok{(}\DecValTok{0}\OperatorTok{,} \DecValTok{0}\OperatorTok{,} \DecValTok{1}\OperatorTok{)}
    \ControlFlowTok{else}
        \VariableTok{delta} \OperatorTok{=} \VariableTok{vector}\OperatorTok{.}\NormalTok{new}\OperatorTok{(}\DecValTok{0}\OperatorTok{,} \DecValTok{1}\OperatorTok{,} \DecValTok{0}\OperatorTok{)}
    \ControlFlowTok{end}

    \VariableTok{delta} \OperatorTok{=} \VariableTok{vector}\OperatorTok{.}\NormalTok{multiply}\OperatorTok{(}\VariableTok{delta}\OperatorTok{,} \VariableTok{dtime}\OperatorTok{)}

    \VariableTok{self}\OperatorTok{.}\VariableTok{object}\OperatorTok{:}\NormalTok{move\_to}\OperatorTok{(}\VariableTok{vector}\OperatorTok{.}\NormalTok{add}\OperatorTok{(}\VariableTok{pos}\OperatorTok{,} \VariableTok{delta}\OperatorTok{))}
\KeywordTok{end}

\KeywordTok{function} \VariableTok{MyEntity}\OperatorTok{:}\NormalTok{on\_punch}\OperatorTok{(}\VariableTok{hitter}\OperatorTok{)}
    \VariableTok{core}\OperatorTok{.}\NormalTok{chat\_send\_player}\OperatorTok{(}
        \VariableTok{hitter}\OperatorTok{:}\NormalTok{get\_player\_name}\OperatorTok{(),}
        \VariableTok{self}\OperatorTok{.}\VariableTok{message}
    \OperatorTok{)}
\KeywordTok{end}
\end{Highlighting}
\end{Shaded}

Wenn du diese Entity erzeugst und benutzt, fällt dir auf, dass die
Nachricht vergessen wird, sobald die Entity inaktiv und wieder aktiv
wird. Das liegt daran, dass die Nachricht nicht gespeichert wird. Statt
einfach die ganze Entity-Tabelle zu speichern, überlässt Luanti dir, was
gespeichert wird. \emph{Staticdata} ist ein String mit allen eigenen
Informationen, die gespeichert werden müssen.

\begin{Shaded}
\begin{Highlighting}[]
\KeywordTok{function} \VariableTok{MyEntity}\OperatorTok{:}\NormalTok{get\_staticdata}\OperatorTok{()}
    \ControlFlowTok{return} \VariableTok{core}\OperatorTok{.}\NormalTok{write\_json}\OperatorTok{(\{}
        \VariableTok{message} \OperatorTok{=} \VariableTok{self}\OperatorTok{.}\VariableTok{message}\OperatorTok{,}
    \OperatorTok{\})}
\KeywordTok{end}

\KeywordTok{function} \VariableTok{MyEntity}\OperatorTok{:}\NormalTok{on\_activate}\OperatorTok{(}\VariableTok{staticdata}\OperatorTok{,} \VariableTok{dtime\_s}\OperatorTok{)}
    \ControlFlowTok{if} \VariableTok{staticdata} \OperatorTok{\textasciitilde{}=} \StringTok{""} \KeywordTok{and} \VariableTok{staticdata} \OperatorTok{\textasciitilde{}=} \KeywordTok{nil} \ControlFlowTok{then}
        \KeywordTok{local} \VariableTok{data} \OperatorTok{=} \VariableTok{core}\OperatorTok{.}\NormalTok{parse\_json}\OperatorTok{(}\VariableTok{staticdata}\OperatorTok{)} \KeywordTok{or} \OperatorTok{\{\}}
        \VariableTok{self}\OperatorTok{:}\NormalTok{set\_message}\OperatorTok{(}\VariableTok{data}\OperatorTok{.}\VariableTok{message}\OperatorTok{)}
    \ControlFlowTok{end}
\KeywordTok{end}
\end{Highlighting}
\end{Shaded}

Luanti darf \texttt{get\_\allowbreak{}staticdata()} beliebig oft und zu
jedem Zeitpunkt aufrufen. Luanti wartet nämlich nicht, bis ein
Mapblock\index{mapblock@Mapblock} inaktiv wird, um ihn zu speichern,
denn das könnte zu Datenverlust führen. Mapblöcke werden ungefähr alle
18 Sekunden gespeichert; in ähnlichen Abständen wird also auch
\texttt{get\_\allowbreak{}staticdata()} aufgerufen.

\texttt{on\_\allowbreak{}activate()}\index{callback@Callback!on_activate@\texttt{on\_\allowbreak{}activate}}
wird dagegen nur aufgerufen, wenn eine Entity aktiv wird -- entweder
weil ihr Mapblock aktiv wird oder weil sie neu erzeugt wird. Staticdata
kann deshalb leer sein.

Zum Schluss registrierst du die Typ-Tabelle mit der passend benannten
Funktion \texttt{register\_\allowbreak{}entity}.

\begin{Shaded}
\begin{Highlighting}[]
\VariableTok{core}\OperatorTok{.}\NormalTok{register\_entity}\OperatorTok{(}\StringTok{"mymod:entity"}\OperatorTok{,} \VariableTok{MyEntity}\OperatorTok{)}
\end{Highlighting}
\end{Shaded}

Ein Mod kann die Entity so erzeugen:

\begin{Shaded}
\begin{Highlighting}[]
\KeywordTok{local} \VariableTok{pos} \OperatorTok{=} \OperatorTok{\{} \VariableTok{x} \OperatorTok{=} \DecValTok{1}\OperatorTok{,} \VariableTok{y} \OperatorTok{=} \DecValTok{2}\OperatorTok{,} \VariableTok{z} \OperatorTok{=} \DecValTok{3} \OperatorTok{\}}
\KeywordTok{local} \VariableTok{obj} \OperatorTok{=} \VariableTok{core}\OperatorTok{.}\NormalTok{add\_entity}\OperatorTok{(}\VariableTok{pos}\OperatorTok{,} \StringTok{"mymod:entity"}\OperatorTok{,} \KeywordTok{nil}\OperatorTok{)}
\end{Highlighting}
\end{Shaded}

Der dritte Parameter sind die anfänglichen Staticdata. Die Nachricht
setzt du über die Methode der Entity-Tabelle:

\begin{Shaded}
\begin{Highlighting}[]
\VariableTok{obj}\OperatorTok{:}\NormalTok{get\_luaentity}\OperatorTok{():}\NormalTok{set\_message}\OperatorTok{(}\StringTok{"Hallo!"}\OperatorTok{)}
\end{Highlighting}
\end{Shaded}

Spieler mit dem
Privileg\index{privileg@Privileg}~(Kapitel~\ref{ch:players/privileges})
\emph{give} können Entities mit einem
Chatbefehl\index{chatbefehl@Chatbefehl}~(Kapitel~\ref{ch:players/chat})
erzeugen:

\begin{Shaded}
\begin{Highlighting}[]
\NormalTok{/spawnentity mymod:entity}
\end{Highlighting}
\end{Shaded}

\section{Lebenspunkte und Schaden}\label{lebenspunkte-und-schaden}

\subsection{Lebenspunkte (HP)}\label{lebenspunkte-hp}

Jedes Objekt hat einen Wert für
Lebenspunkte\index{lebenspunkte@Lebenspunkte} (\emph{health points},
HP), der seine aktuelle Gesundheit angibt. Die maximalen HP legt die
Objekteigenschaft \texttt{hp\_\allowbreak{}max} fest. Fallen die HP auf
0, stirbt das Objekt.

\begin{Shaded}
\begin{Highlighting}[]
\KeywordTok{local} \VariableTok{hp} \OperatorTok{=} \VariableTok{object}\OperatorTok{:}\NormalTok{get\_hp}\OperatorTok{()}
\VariableTok{object}\OperatorTok{:}\NormalTok{set\_hp}\OperatorTok{(}\VariableTok{hp} \OperatorTok{+} \DecValTok{3}\OperatorTok{)}
\end{Highlighting}
\end{Shaded}

\subsection{Schläge, Schadensgruppen und
Rüstungsgruppen}\label{schluxe4ge-schadensgruppen-und-ruxfcstungsgruppen}

Schaden verringert die HP eines Objekts. Ein Objekt kann ein anderes
\emph{schlagen} (\emph{punch}), um ihm Schaden zuzufügen. Ein Schlag
muss kein echter Schlag sein -- es kann auch eine Explosion, ein
Schwerthieb oder etwas anderes sein.

Der Gesamtschaden ergibt sich, indem die
Schadensgruppen\index{schadensgruppe@Schadensgruppe} des Schlags mit der
Verwundbarkeit des Ziels multipliziert werden. Das Ergebnis wird dann
begrenzt, je nachdem, wie kurz der letzte Schlag zurückliegt. Gleich
rechnen wir das an einem Beispiel durch.

Wie die Abbautypen\index{abbautyp@Abbautyp} von
Nodes~(Kapitel~\ref{ch:items/nodes-items-crafting}) können diese
Gruppen\index{gruppe@Gruppe} beliebige Namen haben und müssen nicht
registriert werden. Üblich ist aber, dieselben Gruppennamen wie beim
Abbauen von Nodes zu verwenden.

Wie verwundbar ein Objekt für bestimmte Schadensarten ist, hängt von
seinen \texttt{armor\_\allowbreak{}groups} ab. Trotz des irreführenden
Namens geben \texttt{armor\_\allowbreak{}groups} an, wie viel Prozent
Schaden das Objekt von einer Schadensgruppe erleidet -- nicht, wie viel
es abwehrt. Fehlt eine Schadensgruppe in den Rüstungsgruppen eines
Objekts, ist es gegen diese völlig unverwundbar.

\begin{Shaded}
\begin{Highlighting}[]
\VariableTok{target}\OperatorTok{:}\NormalTok{set\_armor\_groups}\OperatorTok{(\{}
    \VariableTok{fleshy} \OperatorTok{=} \DecValTok{90}\OperatorTok{,}
    \VariableTok{crumbly} \OperatorTok{=} \DecValTok{50}\OperatorTok{,}
\OperatorTok{\})}
\end{Highlighting}
\end{Shaded}

Im Beispiel erleidet das Objekt 90 \% des \texttt{fleshy}-Schadens und
50 \% des \texttt{crumbly}-Schadens.

Schlägt ein Spieler ein Objekt, kommen die Schadensgruppen von dem Item,
das er gerade in der Hand hält. In anderen Fällen legen Mods fest,
welche Schadensgruppen verwendet werden.

\subsection{Beispiel für eine
Schadensberechnung}\label{beispiel-fuxfcr-eine-schadensberechnung}

Schlagen wir das Objekt \texttt{target}:

\begin{Shaded}
\begin{Highlighting}[]
\KeywordTok{local} \VariableTok{tool\_capabilities} \OperatorTok{=} \OperatorTok{\{}
    \VariableTok{full\_punch\_interval} \OperatorTok{=} \DecValTok{0.8}\OperatorTok{,}
    \VariableTok{damage\_groups} \OperatorTok{=} \OperatorTok{\{} \VariableTok{fleshy} \OperatorTok{=} \DecValTok{5}\OperatorTok{,} \VariableTok{choppy} \OperatorTok{=} \DecValTok{10} \OperatorTok{\},}

    \CommentTok{{-}{-} Nur fürs Abbauen von Nodes, aber trotzdem nötig}
    \VariableTok{max\_drop\_level} \OperatorTok{=} \DecValTok{1}\OperatorTok{,}
    \VariableTok{groupcaps} \OperatorTok{=} \OperatorTok{\{}
        \VariableTok{fleshy} \OperatorTok{=} \OperatorTok{\{}
            \VariableTok{times} \OperatorTok{=} \OperatorTok{\{} \OperatorTok{[}\DecValTok{1}\OperatorTok{]} \OperatorTok{=} \DecValTok{2.5}\OperatorTok{,} \OperatorTok{[}\DecValTok{2}\OperatorTok{]} \OperatorTok{=} \DecValTok{1.20}\OperatorTok{,} \OperatorTok{[}\DecValTok{3}\OperatorTok{]} \OperatorTok{=} \DecValTok{0.35} \OperatorTok{\},}
            \VariableTok{uses} \OperatorTok{=} \DecValTok{30}\OperatorTok{,}
            \VariableTok{maxlevel} \OperatorTok{=} \DecValTok{2}\OperatorTok{,}
        \OperatorTok{\},}
    \OperatorTok{\},}
\OperatorTok{\}}

\KeywordTok{local} \VariableTok{time\_since\_last\_punch} \OperatorTok{=}
    \VariableTok{tool\_capabilities}\OperatorTok{.}\VariableTok{full\_punch\_interval}
\VariableTok{target}\OperatorTok{:}\NormalTok{punch}\OperatorTok{(}
    \VariableTok{object}\OperatorTok{,}
    \VariableTok{time\_since\_last\_punch}\OperatorTok{,}
    \VariableTok{tool\_capabilities}
\OperatorTok{)}
\end{Highlighting}
\end{Shaded}

Jetzt berechnen wir den Schaden. Die Schadensgruppen des Schlags sind
\texttt{fleshy=5} und \texttt{choppy=10}; \texttt{target} erleidet 90 \%
des \texttt{fleshy}-Schadens und 0 \% des \texttt{choppy}-Schadens.

Zuerst multiplizieren wir die Schadensgruppen mit der Verwundbarkeit und
addieren die Ergebnisse. Dann multiplizieren wir mit einer Zahl zwischen
0 und 1, die von
\texttt{time\_\allowbreak{}since\_\allowbreak{}last\_\allowbreak{}punch}
abhängt.

\begin{Shaded}
\begin{Highlighting}[]
\NormalTok{t = limit(time\_since\_last\_punch / full\_punch\_interval, 0, 1)}
\NormalTok{  = 1}

\NormalTok{Schaden = (5*90/100 + 10*0/100) * t}
\NormalTok{        = 4.5}
\end{Highlighting}
\end{Shaded}

Da HP ganze Zahlen sind, wird der Schaden auf 5 Punkte gerundet.

\section{Anhängen}\label{anhuxe4ngen}

Angehängte Objekte bewegen sich mit, wenn sich das Elternobjekt -- das
Objekt, an das sie angehängt sind -- bewegt. Ein angehängtes Objekt
heißt Kind des Elternobjekts. Ein Objekt kann beliebig viele Kinder,
aber höchstens ein Elternobjekt haben.

\begin{Shaded}
\begin{Highlighting}[]
\VariableTok{child}\OperatorTok{:}\NormalTok{set\_attach}\OperatorTok{(}\VariableTok{parent}\OperatorTok{,} \VariableTok{bone}\OperatorTok{,} \VariableTok{position}\OperatorTok{,} \VariableTok{rotation}\OperatorTok{)}
\end{Highlighting}
\end{Shaded}

\texttt{get\_\allowbreak{}pos()} gibt immer die globale Position eines
Objekts zurück, egal ob es angehängt ist oder nicht.
\texttt{set\_\allowbreak{}attach} erwartet eine relative Position, aber
nicht so, wie man es erwarten würde: Die Position bezieht sich auf den
Ursprung des Elternobjekts, mit dem Faktor 10 skaliert. \texttt{0,5,0}
ist also einen halben Node über dem Ursprung des Elternobjekts.

\begin{hinweis}{Grad und Bogenmaß}

Die Drehung beim Anhängen wird in Grad angegeben, die Drehung von
Objekten dagegen im Bogenmaß. Achte darauf, in die richtige Einheit
umzurechnen.

\end{hinweis}

Bei animierten 3D-Modellen hängst du die Entity mit dem Argument
\texttt{bone} an einen Knochen. 3D-Animationen basieren auf Skeletten --
einem Netz aus Knochen im Modell, denen jeweils eine Position und
Drehung gegeben werden kann, um das Modell zu verändern, etwa um den Arm
zu bewegen. Das Anhängen an einen Knochen ist nützlich, wenn eine Figur
etwas in der Hand halten soll:

\begin{Shaded}
\begin{Highlighting}[]
\VariableTok{obj}\OperatorTok{:}\NormalTok{set\_attach}\OperatorTok{(}\VariableTok{player}\OperatorTok{,}
    \StringTok{"Arm\_Right"}\OperatorTok{,}           \CommentTok{{-}{-} Standardknochen}
    \OperatorTok{\{}\VariableTok{x}\OperatorTok{=}\DecValTok{0.2}\OperatorTok{,} \VariableTok{y}\OperatorTok{=}\DecValTok{6.5}\OperatorTok{,} \VariableTok{z}\OperatorTok{=}\DecValTok{3}\OperatorTok{\},}   \CommentTok{{-}{-} Standardposition}
    \OperatorTok{\{}\VariableTok{x}\OperatorTok{={-}}\DecValTok{100}\OperatorTok{,} \VariableTok{y}\OperatorTok{=}\DecValTok{225}\OperatorTok{,} \VariableTok{z}\OperatorTok{=}\DecValTok{90}\OperatorTok{\})} \CommentTok{{-}{-} Standarddrehung}
\end{Highlighting}
\end{Shaded}

\section{Jetzt bist du dran}\label{jetzt-bist-du-dran}

\begin{itemize}
\tightlist
\item
  Baue eine Windmühle aus Nodes und einer Entity.
\item
  Erstelle einen Mob deiner Wahl (nur mit der Entity-API und ohne andere
  Mods).
\end{itemize}
